Las medallas secretas se desbloquean con códigos especiales que forman palabras. La página sigue actualizándose y ya son más de 80, así que esta lista puede seguir creciendo, seguramente en el drive estén más actualizadas que acá.

\subsection*{Medallas secretas con código}
\begin{description}[leftmargin=1.4cm, style=nextline]
  \item[3AKJEE] Del creador de Espacio UNAM (el jefazo). Otorga una medalla secreta. --- Atte. Runix.
  \item[T3NAZ7] Se premia nuestra perseverancia (medalla \emph{Tenaz}).
  \item[3Y3E34] Otro código secreto: la UNAMball.
  \item[OLAK4S3] Para el de la llama.
  \item[H3RM4N0] Para el de: Hermandad.
  \item[N0LAT3N] Para el de: ``Esa tú no la tienes''.
  \item[C1UDAD0] Para la de: Ciudadano promedio.
  \item[T0D0B1EN] Para el de: Todo está bien.
  \item[CANSAD0] Para el de: Estoy cansado, jefe.
  \item[N0SERICK] Para el de: No lo sé, Rick.
  \item[TR4NQU1] Para la de: Oye, tranquilo viejo.
  \item[ESTR3LL4] Para la de: Estrellita en la frente. (by Adv)
  \item[CH1NCH3] Por la fama de enchincharnos en la UNAM y por los px de veterinaria- Propuesta por Mariana
  \item[1YAR1] Homenaje a Iyari, la puma de la Facultad de Medicina veterinaria y Zootecnia. Idea e icono de Nahomi. Medalla secreta
  \item[G0Y0] Para Goyo, la mascota del Club Universidad Nacional: el puma que pone a rugir el estadio. Idea e ícono de José Carlos. Medalla secreta
  \item[PR0M3T30] Para el Prometeo de la Facultad de Ciencias. Idea de Emmanuel y Santiago; ícono de Santiago. Medalla secreta. el prometeo llevando el conocimiento hacia nuevos horizontes
  \item[M4ND4R1N] Para Mandarino, el gatito naranja de la Facultad de Química que la comunidad cuidó y esterilizó, y que hoy vive en Canadá. Idea de Zoe
  \item[R3S3RV4] Por la Reserva Ecológica del Pedregal de San Ángel, el espacio vivo dentro de Ciudad Universitaria. Idea de Andrés
  \item[LGBTQ1] Para la comunidad LGBTQ+ de la UNAM y quienes la acompañan: aquí cabe toda la diversidad. Idea de Ulises. Medalla secreta
  \item[4UR4FARM] Idea mía Gael hay una medalla por el top 100 en el mini juego secreto idea mía tmb Y una por proponer una medalla Para quien farmea aura por los pasillos de la UNAM. Idea de Gael
  \item[N0B3LQ] Medalla secreta.
  \item[N0B3LP] Medalla secreta.
  \item[N0B3LL] Investiguen quiénes han sido nuestros premios Nobel!!!
  \item[C0L4M0CH] Para Cola Colamocha, el perro dehCCH que toda su comunidad quiso y recuerda. Idea e ícono de Elian. Medalla secreta. CUPUM4SLa porra más viral de la UNAM. Idea de Ian, Edgar, Aurora(su servidora) y Alan
  \item[4RD1LL4] Para las ardillas de los campus, que nos roban la comida con toda la confianza. Idea de Adrián, Marco (su servilleta) y Leonardo. Medalla secreta
  \item[S4LM0N3L] Comiste en el Pasillo de la Salmonela y sobreviviste para contarlo. Idea de Tomás y Denis. Atte. Dapheriel H4YP4R0 - Para quienes alzan la voz y defienden sus derechos y los de su comunidad. Idea de Andrea M., David, Guadalupe, Leslie y Andrea B. Medalla secreta
  \item[HU3LG4] Para la Muela de la Facultad de Odontología, uno de los rincones más reconocibles de CU. Idea de Sergio y Santiago; ícono de Santiago. Medalla secreta
  \item[M1ZTL1] Homenaje a Miztli, la puma que fue emblema de la UNAM. Idea de Getzemani. Medalla secreta. HUM4N1D4DPara «La Humanidad Hoy», el monumento que es símbolo de la FES Acatlán. Idea e ícono de Hanna. Medalla secreta
  \item[TR3N1NG] Para el tren del edificio principal de la Facultad de Ingeniería. Idea de Israel. Medalla secreta
  \item[0NLIN31159] Para quienes estudian en línea y suben la tarea al último minuto, pero la suben. Idea de Ángel Alexander
  \item[CH1T0] Para Chito, la hormiga mascota del CCH Azcapotzalco. Idea e ícono de Diego Yamil. Medalla secreta
  \item[C0VID19] Para quienes estudiaron en la UNAM durante la pandemia de COVID-19 (2020–2022): clases en línea, cámaras apagadas y aun así seguimos adelante. Idea de Ilse. Medalla secreta
  \item[V14J3R0] Para quienes hacen dosBS-26-279927-4 horas o más de camino a su facultad o sede, sobreviviendo al tráfico en un viaje interdimensional diario: cada día llegan y eso ya es un logro. Idea de Yocelyn y Jhosep. Medalla secreta
  \item[2DEOCT] En memoria del movimiento estudiantil de 1968 y de quienes lo dieron todo: a esa generación le debemos mucho de lo que hoy tenemos. Las cosas han mejorado, pero todavía hay causas por las cuales luchar. «Somos nietos de la Revolución, hijos del 68 y hermanos de los 43.» Idea de Jonathan, Jorge Pavel y Cecilia. Medalla secreta. STCCU Para la estación Universidad, la puerta de entrada a Ciudad Universitaria para miles de pumas cada día. Idea de Luis Enrique y Jorge. Medalla secreta
  \item[M4STR04P4] Una medalla secreta para quien domina el formato APA sin perder la cordura. Idea de Abigail. Medalla Secreta
  \item[JUL1F13RR0] En honor a Julieta Fierro, astrónoma de la UNAM y una de las grandes divulgadoras de la ciencia en México. Idea e ícono de Eder Gael; también la propuso Mariana Jaired. Medalla secreta
  \item[S1XS3V3N] Six… seven. El meme del momento. Solo se gana con su código secreto. Idea de Raúl, Adrián, Isaac
  \item[T4C0S] El desayuno oficial de la comunidad: el de la bici, la canasta y la salsa verde
  \item[PUM4D0L4R] Los pumadólares que te cobran en tu primer día… y que nunca existieron
  \item[P13DR4] Para la piedra más famosa de la FES Acatlán y su leyenda: si te tomas foto con ella, no te titulas. Idea de Alyzon, Jair, Eibraham y 9 personas más. Medalla secreta
  \item[3SP1G4] Para «La Espiga», la escultura que aparece en el escudo de la FES Zaragoza. Idea de Sergio, Hernán, César y 8 personas más. Medalla secreta
  \item[UN4MOR] Para quienes encontraron el amor en la UNAM: el campus también es lugar de encuentros. Idea de Elian, Irene, Denis y 7 personas más. Medalla secreta
  \item[BU1T0FQ] Para el búho que recibe a la comunidad de la Facultad de Química. Idea de Marian, Raúl Alberto, Noel Omar y 6 personas más. Medalla secreta
  \item[S0RJU4N4] Para el monumento a Sor Juana Inés de la Cruz de la FES Acatlán, que distingue a toda su comunidad. Idea de Andrea y Johanna. Medalla secreta. [tu propio codigo] Para la medalla secreta de amor propio
  \item[GR4N1T0] Para quien pone un poco de sí para hacer de este mundo un lugar mejor. Medalla secreta
  \item[G1S3S] Para «Los Gises», la escultura de Helen Escobedo que identifica a la FES Iztacala. Idea de Cynthia, Joshua, Alyzon y 14 personas más; ícono de Edgar David. Medalla secreta
  \item[4XOLOT] Para los ajolotes, tan queridos en la FAD de Xochimilco. Idea de Nathalia. Medalla secreta
  \item[AGU1LA] Para el águila del pilar de la UNAM, la facultad de Derecho. Idea de Zyanya, Frida, Javier y 5 personas más. Medalla secreta
  \item[T4K1S] Los Takis de Zaragoza Para «El abrazo del maíz», la escultura del Campus I de la FES Zaragoza que la comunidad bautizó como los Takis. Idea de Laura Fabiola, Lluvia, Leonardo y 4 personas más. Medalla secreta
  \item[BUF4L1N] Para Bufalín, la mascota de la Facultad de Contaduría y Administración. Idea de Luis Enrique, Mayte, Karla Angélica y 6 personas más. Medalla secreta
  \item[4L1N3] Serie «Mujeres que hicieron la UNAM»: Aline Schunemann de Aluja, médica veterinaria y primera mujer profesora emérita de la UNAM. Idea de María Concepción y Karla Beatriz. Medalla secreta
  \item[D3LF1N] Para el delfín, el animal que representa a la ENP 7 «Ezequiel A. Chávez». Idea de Emmanuel, Sofía Lucrecia, Claudia e Itzel. Medalla secreta
  \item[T0T3M] Para el Atlante de Tula de la ENP 9, al que la comunidad le deja ofrendas y monedas antes de los exámenes. Idea de Fernanda Camila, José Emilio, Leonardo y 5 personas más. Medalla secreta
  \item[HUG0L] Para Hugo Sánchez: estudió Odontología en la UNAM y en 2026 se cumplen 50 años de su debut con Pumas. Idea de Ana Isabella, Ángel Benjamín, Braulio y 6 personas más. Medalla secreta
  \item[4RB0LJUM3X] Para el árbol de la Facultad de Química que toda su comunidad conoce como «el árbol Jumex». Idea de Abril, Elsa Daniela, Gonzalo y 4 personas más. Medalla secreta
  \item[M1L3T0] En memoria de Mario Lenguas de Mileto, el perrito del CCH Naucalpan que acompañó a muchas generaciones y hoy tiene un mural en su biblioteca. Idea de María Jasmín, Arcelia, Andrea y 6 personas más. Medalla secreta
  \item[V4MP1RO] Medalla secreta.
  \item[TUTS1] Para «MéXXico, el siglo lleva tu nombre», la escultura de la Facultad de Arquitectura que la comunidad llama la Tutsi Pop. Idea de Sergio, Oswald, Axel Emmanuel y 7 personas más. Medalla secreta
  \item[B1G0T3S] Para Los Bigotes, la escultura de CU junto al Metro Universidad y la parada del Pumabús. Idea de Sergio, Raúl, Diego Antonio y 4 personas más. Medalla secreta
  \item[B1NN104X] Para Binni Sicarú, la nutria mascota de la ENES Oaxaca. Idea d un e Óscar Isidro, Excel Zuriel y Alfredo Daniel. Medalla secreta
  \item[1SL4SCU] Hay lugares que se vuelven parte de nuestra historia: Las Islas son el corazón de Ciudad Universitaria. Idea de Celeste. Medalla secreta
  \item[C0L1BR1] Para el colibrí de Trabajo Social, que dejó de ser Escuela y ahora es Facultad: un momento histórico. A su comunidad le llega sola. Idea de Karen, Karina, Daniela y 5 personas más. Medalla secreta
  \item[TUZ4S0R] Para las tuzas que viven en los pastos del CCH Oriente. Idea de Santiago, Naomi, Rafael y 2 personas más. Medalla secreta
  \item[4ZUL0R0] Por los colores que nos identifican en todo el mundo, idea de Adri. Medalla Secreta
  \item[G0Y4] Para Goya, la puma del equipo femenil de Pumas y compañera de Goyo. Idea de Giselle, Priscilla, Pedro y 7 personas más; ícono de Pedro. Medalla secreta
  \item[ESPUM1T4] Para el jabón que inventaron en la UNAM: hasta la espuma puede ser ciencia universitaria. Idea de Dulce Liliana. Medalla secreta
  \item[IBM650] Para la IBM 650, la primera computadora de México, que llegó a la UNAM en 1958. Idea de Elizabeth Briseida. Medalla secreta
  \item[M0NT0Y4] Para Matilde Montoya, la primera médica de México: se tituló en la Escuela Nacional de Medicina en 1887. Serie «Mujeres que hicieron la UNAM». Idea de Ángel Manuel y Audrie Annel. Medalla secreta
  \item[3PP3NSFM] Para «La vida, la muerte, el mestizaje y los cuatro elementos», el mural de Francisco Eppens en la Facultad de Medicina. Idea de Diego Alejandro, Jessica Alexa, Rodrigo y 2 personas más. Medalla secreta
  \item[PR0F3S] Para las y los profes y las personas trabajadoras que sostienen la UNAM todos los días. Idea de Héctor Erick, Dalthon y Carol Lizeth. Medalla secreta
  \item[H4RU] Para Haru, el gatito del CCH Vallejo que tiene su propio mural en el plantel. Idea de Maricruz Martínez y Silvia Pozos. Serie «Mascotas de mi escuela». Medalla secreta
  \item[0C3L0T3] Para el Ocelote, la mascota que lleva con orgullo los colores y los equipos de la FES Zaragoza. Idea de Juan Hernández y Silvia Pozos. Serie «Mascotas de mi escuela». Medalla secreta
  \item[3SP3R4NZ4] Para quienes siguen adelante con valentía en sus propias batallas: siempre hay una luz al final del camino y una comunidad que camina contigo. Idea de Omar Ortega, Jhosep Silva y Dalia Maldonado
  \item[S0NRIS4] Para Ris, el león dientes de sable que le saca una sonrisa a toda la Facultad de Odontología. Idea de Jesús Núñez. Serie «Mascotas de mi escuela». Medalla secreta
  \item[BUH0] Para el Búho, mascota de la Facultad de Derecho: sabiduría y desvelos de estudio. Idea de Gojan Neri. Serie «Mascotas de mi escuela». Medalla secreta
  \item[4J0L0T3] Para el ajolote, mascota del CCH Vallejo. Idea de José Contreras. Serie «Mascotas de mi escuela». Medalla secreta
  \item[L30P4RD0] Para los Leopardos, la mascota que ruge por la Prepa 8 «Miguel E. Schulz». Idea de Eduardo Palacio y Emiliano Zamudio. Serie «Mascotas de mi escuela». Medalla secreta. Para los Coyotes, emblema y equipos de la Prepa 6 «Antonio Caso» en Coyoacán. Idea de Raúl Coronado, Elías Macotela, Amos Vega y más. Serie «Mascotas de mi escuela». Medalla secreta
  \item[0SC4R] Para Alfonso Cuarón y Emmanuel Lubezki, que estudiaron cine en la UNAM y ganaron el Óscar. Idea de Dylan, Denis, Frida Sabina y Letizia. Medalla secreta
  \item[DG1R3] Para la comunidad del Sistema Incorporado: preparatorias, CCH y licenciaturas incorporadas que también son parte de la UNAM. Idea de Mariana Celeste. Medalla secreta
  \item[B3T1T0] Para Betito, el capibara mascota cultural de la Facultad de Ingeniería. Idea de Emanuel Escobar y Andrea Rosales. Serie «Mascotas de mi escuela». Medalla secreta
  \item[B0PUM3] Para el buque oceanográfico El Puma de la UNAM, que estudia los mares de México. Idea de Salvador y Ximena. Medalla secreta
  \item[CH4Y0] Para Rosario Castellanos, escritora que estudió en la Facultad de Filosofía y Letras. Serie «Mujeres que hicieron la UNAM». Idea de Alejandra
  \item[H3L14] Serie «Mujeres que hicieron la UNAM»: Helia Bravo Hollis, primera bióloga titulada en México y gran estudiosa de los cactus. Idea de Víctor Ramón, Mariana Jaired, Yelena y Brenda. H4LCON /
  \item[PR3P1N] Para Prepín y los Halcones, el orgullo de la Prepa 9 «Pedro de Alba». Idea de Jonathan Álvarez, Katia Chávez, Eduardo Palacio y más. Serie «Mascotas de mi escuela»
  \item[CH4T4] Para Chata, la perrita más querida de la FES Zaragoza Campus 2, compañera fiel de toda su comunidad. Idea de Fátima Cruz y Dana García. Serie «Mascotas de mi escuela». Medalla secreta
  \item[V41N1LL4] Para Vainilla, la husky consentida de la ENES Morelia. Idea de Montserrat Miranda y Valeria Olmos. Serie «Mascotas de mi escuela». Medalla secreta. En total, son muchas medallas secretas mis estimados (al menos eso me dijo GoIA, además de que siguen actualizando la página). Atte. Diego Terán Ahora son 200 JAJAJA no paran de agregar más. Una disculpa JAJAJA La UNAM dándose cuenta que no fue buena idea aceptar solicitudes (no van a parar) se acabarán las ideas en algún momento ntp (o se dan por vencidos, 1 de 2) LAS MEDALLAS NO PARAN DE SPAWNEAR WTF JSJD consecuencias de dejarnos sugerirlas. Ya me cansé de agregar los códigos secretos, mañana amaneceremos con 100 JAJA. Buenas noches, mis estimados adictos a los lazos
  \item[T0RR3S] Homenaje a las torres de la FES Aragón
  \item[TL4CU4] Medalla de tlacuache. Medalla propuesta por Ximena Valencia
  \item[P0K3B0L4] La Leonardita del anexo de Ingeniería, a la que todo el campus conoce como la Pokébola. Idea de Alexis e Israel; ícono de Israel. Medalla secreta. Atte. Dapheriel
  \item[CUPUM4S] La porra más viral de la UNAM. Idea de Ian, Edgar, Aurora(su servidora) y Alan.
  \item[H4YP4R0] Medalla secreta.
  \item[4BR3C0CA] Para la Muela de la Facultad de Odontología, uno de los rincones más reconocibles de CU. Idea de Sergio y Santiago; ícono de Santiago. Medalla secreta.
  \item[HUM4N1D4D] Para «La Humanidad Hoy», el monumento que es símbolo de la FES Acatlán. Idea e ícono de Hanna. Medalla secreta.
  \item[C0Y0T3] Para los Coyotes, emblema y equipos de la Prepa 6 «Antonio Caso» en Coyoacán. Idea de Raúl Coronado, Elías Macotela, Amos Vega y más. Serie «Mascotas de mi escuela». Medalla secreta.
  \item[H4LCON] Medalla secreta.
  \item[1ZT4YC4L4] Para Izta y Cala, el puma y la puma que acompañan y motivan a toda la comunidad de la FES Iztacala. Idea de Andrea Cario, José Millán y Axel García. Serie «Mascotas de mi escuela». Medalla secreta.
  \item[H0L4MD0] Medalla secreta.
  \item[H0L4WL] Medalla secreta.
  \item[GR4C1AZ] Para las mascotas que se desvelan a tu lado cuando estudias o preparas una exposición: también se merecen su medalla. Idea de Mariana Adrián. Medalla secreta.
  \item[L0M1T0] Para Lomitos FAD, el grupo de estudiantes de la Facultad de Artes y Diseño que rescata y cuida a los perritos y gatitos de su facultad
  \item[G4R1N] Para Manuela Garín Pinillos, de las primeras mujeres en estudiar matemáticas en la UNAM y pionera de esa ciencia en México. Serie «Mujeres que hicieron la UNAM». Idea de Joaquín Omar. Medalla secreta.
  \item[M4R1ASU] Serie «Mujeres que hicieron la UNAM»: María Asunción Sandoval, primera mujer abogada titulada en México (1898). Idea de Yari Natalia y Ana Esperanza. Medalla secreta.
  \item[1F1G3N1A] Serie «Mujeres que hicieron la UNAM»: Ifigenia Martínez, economista y primera mujer en dirigir la Escuela Nacional de Economía, hoy Facultad de Economía. Idea de Daniela. Medalla secreta.
  \item[P1K4CHU] Para el «Pikachu», como la comunidad de la FES Acatlán llama a su monumento más famoso. Idea de Iván, Eibraham, Johan y 6 personas más. Medalla secreta.
  \item[YAPORF4VOR] Para quienes ya no pueden con tantas medallas nuevas…y aún así siguen juntándose. Idea de Diego, Sara y Oziel. Solo se gana con su código secreto.
  \item[CH0RN3] Para Margarita Chorné y Salazar, la primera dentista titulada de México (1886) y una de las primeras mujeres con título profesional en América Latina. Serie «Mujeres que hicieron la UNAM». Idea de Daniel Flores. Medalla secreta.
  \item[W4FFL3S] Para Waffles, el perro que todo el CCH Oriente conoce. Idea de Brian Lucas. Serie «Mascotas de mi escuela». Medalla secreta.
  \item[TURN0S] Para quienes comparten salones, pasillos y profes pero casi nunca se ven: el turno de la mañana y el de la tarde son una sola comunidad. Idea de Sarahí Guerrero. Medalla secreta.
  \item[VIRUS] Se contagia: alguien con la medalla te lo debe pasar al agregarte. Portadores conocidos: \codigos{5UQAMY} (lo reporta el Drive), \codigos{8KTVEE} (comprobado) y \codigos{V569QT}.
  \item[J41D4R] Para Alejandra Jaidar, una de las primeras físicas de México, egresada de la Facultad de Ciencias y pionera de la física nuclear en la UNAM
  \item[C4CHORR0] Medalla secreta.
  \item[3R4SM0] Medalla secreta.
  \item[ZOSC9R] Encontraste el código secreto de El Apagón: en su primer día, a Espacio UNAM se le fue la luz… y aun así encendió.
  \item[H1L4R10] Para Hilario, el perrito que se volvió mascota de la Prepa 3 «Justo Sierra». Idea de Ximena Hernández. Serie «Mascotas de mi escuela». Medalla secreta.
  \item[CU4TR3R0S] La porra de la Prepa 4 «Vidal Castañeda y Nájera», la casa de los Cuatreros de Tacubaya. Idea de Ángela Marina Cordero. Serie «Porras». Medalla secreta.
  \item[SU4D3R0] Para Suadero, el perrito de la FES Cuautitlán. Idea de Alan Salvador Laguna. Serie «Mascotas de mi escuela». Medalla secreta.
  \item[OAX4C4] (ya esta bn)
  \item[JUG4D4] Para esa movida maestra que nadie vio venir: hasta un peón puede llegar al otro lado del tablero y coronarse. Medalla secreta.
  \item[4T3N34] Para Atenea, la diosa de la sabiduría que es el emblema de la Facultad de Filosofía y Letras. Idea de Jaisha, Jesús Horacio, Alejandra y 5 personas más. Serie «Mascotas de mi escuela». Medalla secreta.
  \item[M0NJ3] Para el monje, la mascota de la Prepa 3 «Justo Sierra». Idea de Marian y José Antonio. Serie «Mascotas de mi escuela». Medalla secreta.
  \item[0G0RM4N] Para Juan O'Gorman, autor del mosaico de piedra que cubre la Biblioteca Central. Idea de César Adrián y Adriana. Serie «Personajes de la UNAM». Medalla secreta.
  \item[P4J4R1T0] Para la cara de puma del escudo de los Pumas, que diseñó Manuel «El Pajarito» Andrade. Idea de Aldo Patricio, Adriana y Víctor René. Medalla secreta.
  \item[Y4C4FC4] Para Yacatecuhtli, el dios del comercio que es ícono de la Facultad de Contaduría y Administración. Idea de Edson, Rodrigo, Elisa y Myriam. Serie «Mascotas de mi escuela». Medalla secreta.
  \item[J0RG3C4MP0] Para Jorge Campos, el portero de los uniformes de colores, emblema de los Pumas. Idea de Bogart y Fernanda. Medalla secreta
  \item[4V3F3N1X] Para el ave fénix, el animal que representa a la Facultad de Psicología. Idea de Montserrat y Cecilia. Serie «Mascotas de mi escuela». Medalla secreta.
  \item[P4NT3R4S] Para las Panteras, la mascota de la Facultad de Ciencias Políticas y Sociales y de sus equipos. Idea de Andrés, Isaac y Gustavo. Serie «Mascotas de mi escuela». Medalla secreta.
  \item[C4S4N0V4] Para Pablo González Casanova, el rector que impulsó la creación del CCH y del Sistema Universidad Abierta. Idea de Andrés, Alejandro y Miranda. Serie «Personajes de la UNAM». Medalla secreta.
  \item[L4OCT4V4] Para la afición de los Pumas, en las buenas y en las malas, esperando juntos el octavo título. Idea de Jesús Eliuth y Denis. Medalla secreta.
  \item[B4RR0S] Para Javier Barros Sierra, el rector que defendió la autonomía universitaria en 1968. Idea de Miguel Ángel, José Iyaquibalam y Osvaldo. Serie «Personajes de la UNAM». Medalla secreta.
  \item[L3NGU4S] Para Lengüitas, la mascota del CCH Naucalpan. Idea de María Ximena. Serie «Mascotas de mi escuela». Medalla secreta.
  \item[C0C0DR1L0] Para el cocodrilo de la Prepa 5 «José Vasconcelos». Idea de Óscar Antonio y Ricardo. Serie «Mascotas de mi escuela». Medalla secreta
  \item[T3HU4C4N] La porra de la Prepa 6 «Antonio Caso», de Coyoacán. Idea de Bruno Alvarado. Serie «Porras». Medalla secreta.
  \item[M3ND1Z4B4L] el de la medalla «Concepción Mendizábal»
  \item[K4R1M] Para Karim, el gatito que adoptó la comunidad de la Escuela Nacional de Ciencias Forenses.
  \item[0ANSPM] Para el Observatorio Astronómico Nacional de San Pedro Mártir, en Baja California, donde la UNAM mira el cielo. Idea de Jonathan Uriel, Juan Carlos y Adrián. Serie «Ciencia UNAM». Medalla secreta.
  \item[R1V3R4] Para Diego Rivera, autor del relieve de piedras de colores del Estadio Olímpico Universitario. Idea de Miguel Ángel y Adriana. Serie «Personajes de la UNAM». Medalla secreta.
  \item[R4Y1TAS] Para Rayitas, el perrito que siempre tuvo un lugar en el corazón del CCH Vallejo. Idea de Enrique Moreno y José Gómez. Serie «Mascotas de mi escuela». Medalla secreta.
  \item[BV1TR3S] Para el Buitre, mascota de la Facultad de Medicina y de sus equipos. Idea de Jeremy Hernández. Serie «Mascotas de mi escuela». Medalla secreta.
  \item[L4PR1M3R4] Para los siete títulos de liga de los Pumas, desde el primero en 1977 hasta el de 2011. Idea de Adán, Jesús Alberto y Daniel. Medalla secreta.
  \item[TH3M1S] Para Themis, la diosa de la justicia que acompaña a la comunidad de la Facultad de Derecho. Idea de Elena y Casandra. Serie «Mascotas de mi escuela». Medalla secreta.
  \item[4LCUB13RR3] Para Miguel Alcubierre, físico de la UNAM famoso en el mundo por su propuesta de un motor de curvatura. Idea de Penélope. Serie «Personajes de la UNAM». Medalla secreta.
  \item[UNAM1551] La Universidad cumple 475 años (1551–2026) y tú estuviste ahí: encontraste el código escondido en Espacio. Idea de Luis Jair, Sergio, Ariana y 9 personas más. Medalla secreta.
  \item[GR4MMYS] Para Gabriela Ortiz Torres, compositora formada en la Facultad de Música de la UNAM y ganadora de tres premios Grammy: la música de la UNAM en el escenario más grande del mundo. Idea de Edurne, Jared y Hana. Medalla secreta.
  \item[MP0NC3] Para Manuel M. Ponce, el compositor que dirigió la escuela de música de la Universidad en 1945. Idea de Christofer. Serie «Personajes de la UNAM». Medalla secreta.
  \item[PUM4M0R] Medalla secreta.
  \item[N0STR3S] Medalla secreta.
  \item[C4C0M1X] Para el cacomixtle, el vecino nocturno de la Reserva Ecológica del Pedregal. Idea de Brandon, Emmanuel, Diego y 5 personas más. Medalla secreta.
  \item[0FUN4M] Para la Orquesta Filarmónica de la UNAM, la música que llena la Sala Nezahualcóyotl. Idea de José Emilio, Jazmín y Edwin Fabián. Medalla secreta.
  \item[V4C4CO4P4] Alfalfa, vacas y caca… viva la Facultad de Coapa!» Para la porra y la mascota más querida de la Prepa 5. Idea de Sofía Santiago. Serie «Mascotas de mi escuela». Medalla secreta.
  \item[1999UNAM] En memoria de la huelga estudiantil de 1999 por una universidad pública y gratuita. Idea de José, Fátima y César Fidel. Medalla secreta.
  \item[H1K4RU] Para Hikaru, el panda mascota del Departamento de Psiquiatría y Salud Mental de la Facultad de Medicina. Idea de Xavier. Serie «Mascotas de mi escuela». Medalla secreta.
  \item[JUST0S] Para Justo Sierra, que fundó la Universidad Nacional en 1910. Idea de Johan, Dylan y Saúl. Serie «Personajes de la UNAM». Medalla secreta.
  \item[4STR0N4UT4] Para Rodolfo Neri Vela, egresado de Ingeniería de la UNAM y primer astronauta mexicano (1985). Idea de Jonathan y Dylan. Ah Medalla secreta
  \item[3GU1AL1S] Para Joaquín Eguía Lis, el primer rector de la Universidad Nacional (1910). Idea de Dylan. Serie «Personajes de la UNAM». Medalla secreta.
  \item[G4C3T0] Para Gaceto, el perrito más querido del CCH Vallejo. Idea de Jhoselyn Rosales. Serie «Mascotas de mi escuela». Medalla secreta.
  \item[P4NCH0] Para Pancho, el perrito panzón que todo el CCH Naucalpan conoce. Idea de Diego Valdespino. Serie «Mascotas de mi escuela». Medalla secreta.
  \item[TR1D1L0S4] Para Heberto Castillo, ingeniero y profesor de la UNAM, creador de la tridilosa. Idea de David. Serie «Personajes de la UNAM». Medalla secreta.
  \item[S4LMY] Para Salmy, el perrito rescatado que pasea con su ropita por la Facultad de Química y promueve su proyecto de jabones de aceite reciclado. Idea de Alberto Ramírez, Victoria Romero y Carol Sánchez. Serie «Mascotas de mi escuela». Medalla secreta.
  \item[L30F4RD0] Para LeoFArdo, el leopardo mascota de la Facultad de Arquitectura. Idea de Sophia, Joel y Santiago. Serie «Mascotas de mi escuela». Medalla secreta.
  \item[TUMB4S5] Foto de las Tumbas de la Prepa 5, uno de los rincones más famosos (y misteriosos) de la ENP 5. ¿Ya egresaste? También se gana con el código que guarda la comunidad de la Prepa 5. Idea de Xali; el código, de Alejandro Josué, Xali, Celeste y Emiliano
  \item[BR0WN13] Para Brownie y todos los perritos que acompañan a la comunidad de la ENES León. Idea de Camila López, Miroslava Padilla y Emiliano Barreto. Serie «Mascotas de mi escuela». Medalla secreta.
  \item[R4D10UNAM] Para Radio UNAM, que desde 1937 lleva la voz de la Universidad a todo el país. Idea de Denis, Emmanuel André, Jeshua, Zuri y Leonardo Matthew. Medalla secreta.
  \item[M1PR3C10S0] Para quienes quieren todas las medallas secretas y las cuidan como un anillo: mi precioso. Idea de Rodrigo. Solo se gana con su código secreto
  \item[F0M0PUMA] Empezaste a llenar tu Pasaporte porque todo mundo lo estaba llenando y no te lo ibas a perder. Idea de Rodrigo. Solo se gana con su código secreto.
  \item[V4C4ESF3R4] Supongamos una vaca esférica en el vacío... El chiste clásico de Física, para quienes simplifican el problema hasta que sale. Idea de Metztli. Medalla secreta
  \item[C0NQU3S0] La pregunta más famosa de CCH azcapo. ¿Conqueso tus papas?
  \item[G4B1N1T0] Para Gabinito y los perritos de la Prepa 1 «Gabino Barreda». Idea de Giann Enriachetti y Marlen Nava. Serie «Mascotas de mi escuela». Medalla secreta.
  \item[N0SIST3MA] Medalla secreta.
\end{description}

\subsection*{Medallas automáticas (sin código)}
\begin{description}[leftmargin=3.6cm, style=nextline]
  \item[Persona Famosa] Se obtiene al llegar a los 1500 lazos.
  \item[Bug Killer] Se obtiene al reportar un bug.
\end{description}
